% ==================================================================
% id:       prob-009
% topic:    מבוא להסתברות
% subtopic: עקרון ההכלה וההפרדה
% solution: full
% ==================================================================
במפעל טקסטיל יש שתי מכונות לחיתוך בד. ההסתברות שהמכונה הראשונה תתקלקל ביום מסוים היא $0.1$, ההסתברות שהשנייה תתקלקל היא $0.05$, וההסתברות שלפחות אחת מהן תתקלקל היא $0.12$.
\begin{enumerate}
\item מה ההסתברות ששתי המכונות יתקלקלו?
\begin{solbox}
נסמן $A$ = „הראשונה מתקלקלת”, $B$ = „השנייה מתקלקלת”. לפי הכלה והפרדה, $\PP(A\cap B)=\PP(A)+\PP(B)-\PP(A\cup B)=0.1+0.05-0.12=\ans{0.03}$.
\end{solbox}
\item מה ההסתברות שבדיוק מכונה אחת תתקלקל?
\begin{solbox}
„בדיוק אחת” הוא „לפחות אחת” בלי „שתיהן”: $\PP(A\cup B)-\PP(A\cap B)=0.12-0.03=\ans{0.09}$.
\end{solbox}
\item מה ההסתברות שלפחות מכונה אחת תפעל?
\begin{solbox}
זה המשלים של „שתיהן מתקלקלות”: $1-0.03=\ans{0.97}$.
\end{solbox}
\item מה ההסתברות שרק המכונה הראשונה תתקלקל?
\begin{solbox}
$\PP(A\cap B^C)=\PP(A)-\PP(A\cap B)=0.1-0.03=\ans{0.07}$.
\end{solbox}
\end{enumerate}
